
\subsection{定义 $A_\infty$ 双模同态的障碍 I}
\label{subsec:5-2-5}

本节余下的 5.2.5--5.2.7 小节
将用于证明定理 5.2.35。
为此，需要建立 $A_\infty$ 双模同态的
障碍理论；本小节先说明所需概念。

从非滤过情形开始。
对给定的两个正整数 $K,L$，
先定义 $A_{K,L}$ 双模\liftIndex{I189}
及其同态。
令 $\overline C_1$ 为 $A_K$ 代数，
$\overline C_0$ 为 $A_L$ 代数。
如定义 4.4.6，有余导子
\translatornote{原书实体页 300
将下列两个 $\widehat{\overline d}$
称作 “coalgebra homomorphisms”；
但定义 4.4.6 前明定它们
是由运算延拓的次数 $+1$ 余导子，
而非余代数同态。见疑误 A-0282。}
\[
\begin{aligned}
\widehat{\overline d}^{(1)}_{1\cdots K}:
B_{1\cdots K}\overline C_1[1]
&\longrightarrow
B_{1\cdots K}\overline C_1[1],\\
\widehat{\overline d}^{(0)}_{1\cdots L}:
B_{1\cdots L}\overline C_0[1]
&\longrightarrow
B_{1\cdots L}\overline C_0[1].
\end{aligned}
\]
对分次 $R$-模 $\overline D$，
考虑一列次数为 $+1$ 的运算
\[
\overline{\mathfrak n}_{k_1,k_0}:
B_{k_1}\overline C_1[1]
\otimes\overline D[1]
\otimes B_{k_0}\overline C_0[1]
\longrightarrow\overline D[1],
\qquad k_1\leq K,\quad k_0\leq L.
\]
把它们延拓为双余导子
\[
\begin{aligned}
\widehat{\overline d}_{1\cdots K,1\cdots L}:
&B_{1\cdots K}\overline C_1[1]
\otimes\overline D[1]
\otimes B_{1\cdots L}\overline C_0[1]\\
&\longrightarrow
B_{1\cdots K}\overline C_1[1]
\otimes\overline D[1]
\otimes B_{1\cdots L}\overline C_0[1].
\end{aligned}
\]
其定义是
\[
\begin{aligned}
&\widehat{\overline d}_{1\cdots K,1\cdots L}
(\mathbf x_1\otimes y\otimes\mathbf x_0)\\
&\quad=
\sum_{a,b}(-1)^{\deg'\mathbf x'_{1,a}}
\mathbf x'_{1,a}\otimes
\overline{\mathfrak n}
(\mathbf x''_{1,a}\otimes y\otimes
 \mathbf x'_{0,b})
\otimes\mathbf x''_{0,b}\\
&\qquad\quad+
\widehat{\overline d}^{(1)}_{1\cdots K}
\mathbf x_1\otimes y\otimes\mathbf x_0\\
&\qquad\quad+
(-1)^{\deg'\mathbf x_1+\deg' y}
\mathbf x_1\otimes y\otimes
\widehat{\overline d}^{(0)}_{1\cdots L}
\mathbf x_0.
\end{aligned}
\]
这里按约定 (Conv. 3)，即式 (1.36)，
使用张量的分解记号
$\mathbf x_i
=\sum\mathbf x'_{i,a}\otimes
\mathbf x''_{i,a}$。

\begin{definition}{5.2.44}\label{def:5-2-44}
\begin{enumerate}[label=\textup{(\arabic*)},leftmargin=2.8em]
\item 若
\[
\widehat{\overline d}_{1\cdots K,1\cdots L}
\circ
\widehat{\overline d}_{1\cdots K,1\cdots L}
=0,
\]
则称
$(\overline D,\overline{\mathfrak n})$
为 $(\overline C_1,\overline C_0)$
上的 $A_{K,L}$ 双模。

再设 $\overline C'_1$ 为 $A_K$ 代数、
$\overline C'_0$ 为 $A_L$ 代数，
且 $\overline D'$ 为
$(\overline C'_1,\overline C'_0)$
上的 $A_{K,L}$ 双模。
设
$\overline{\mathfrak f}^{(1)}:
\overline C_1\to\overline C'_1$
为 $A_K$ 同态，
$\overline{\mathfrak f}^{(0)}:
\overline C_0\to\overline C'_0$
为 $A_L$ 同态。
考虑映射
\[
\overline\varphi_{k_1,k_0}:
B_{k_1}\overline C_1[1]
\otimes\overline D[1]
\otimes B_{k_0}\overline C_0[1]
\longrightarrow\overline D'[1],
\quad k_1\leq K,\quad k_0\leq L.
\]
它们连同
$\overline{\mathfrak f}^{(i)}$
诱导
\[
\begin{aligned}
\widehat{\overline\varphi}_{k_1,k_0}:
&B_{1\cdots k_1}\overline C_1[1]
\otimes\overline D[1]
\otimes B_{1\cdots k_0}\overline C_0[1]\\
&\longrightarrow
B_{1\cdots k_1}\overline C'_1[1]
\otimes\overline D'[1]
\otimes B_{1\cdots k_0}\overline C'_0[1].
\end{aligned}
\]
\item 若所有
$\widehat{\overline\varphi}_{k_1,k_0}$
（$k_1\leq K,\ k_0\leq L$）
都是上链同态，
则称 $\overline\varphi_{k_1,k_0}$
组成位于
$(\overline{\mathfrak f}^{(1)},
  \overline{\mathfrak f}^{(0)})$
之上的 $A_{K,L}$ 双模同态
\liftIndex{I190}。
\end{enumerate}
\end{definition}

注意，若 $K\leq K'$、$L\leq L'$，
则 $A_{K',L'}$ 结构诱导 $A_{K,L}$ 结构。
下面描述把 $A_{K,L}$ 双模同态
延拓为 $A_{K+1,L}$ 或
$A_{K,L+1}$ 双模同态的障碍\liftIndex{I350}。
考虑 $R$-模
\[
\operatorname{Hom}\bigl(
B_{K+1}\overline C_1[1]
\otimes\overline D[1]
\otimes B_L\overline C_0[1],
\overline D'[1]\bigr),
\]
并在其上定义上边界算子
\[
\begin{aligned}
\delta_1:
&\operatorname{Hom}\bigl(
B_{K+1}\overline C_1[1]
\otimes\overline D[1]
\otimes B_L\overline C_0[1],
\overline D'[1]\bigr)\\
&\longrightarrow
\operatorname{Hom}\bigl(
B_{K+1}\overline C_1[1]
\otimes\overline D[1]
\otimes B_L\overline C_0[1],
\overline D'[1]\bigr),
\end{aligned}
\]
其作用为
\begin{equation}
\begin{aligned}
\delta_1(\mathcal T)
={}&\overline{\mathfrak n}'_{0,0}
\circ\mathcal T\\
&-(-1)^{\deg\mathcal T}
\mathcal T\circ\bigl(
\widehat{\overline{\mathfrak m}}_1
  \widehat\otimes 1\widehat\otimes1
+1\widehat\otimes
  \overline{\mathfrak n}_{0,0}
  \widehat\otimes1
+1\widehat\otimes1
  \widehat\otimes
  \widehat{\overline{\mathfrak m}}_1
\bigr).
\end{aligned}
\tag{5.2.45}\label{eq:5-2-45}
\end{equation}
式 (5.2.45) 的第一项取靶双模微分
\translatornote{原书实体页 301
第一项写
$\overline{\mathfrak n}_{0,0}
\circ\mathcal T$，
但 $\mathcal T$ 的靶是
$\overline D'[1]$；
作用于它的微分应为
$\overline{\mathfrak n}'_{0,0}$。
括号内作用于源双模
$\overline D[1]$ 的
$\overline{\mathfrak n}_{0,0}$
则保持不变。见疑误 A-0283。}。
这里
\[
\widehat{\overline{\mathfrak m}}_1:
B\overline C_i[1]
\longrightarrow
B\overline C_i[1]
\]
是由 $\overline C_i$ 上的
$\overline{\mathfrak m}_1$
诱导的余导子（见式 (7.1)）。
为简化记号，在不会混淆时，
对 $i=0,1$ 的两侧代数
共用
$\widehat{\overline{\mathfrak m}}_1$
和 $\overline{\mathfrak m}$ 两个记号。

\begin{lemma}{5.2.46}\label{lem:5-2-46}
若
$\overline\varphi:
\overline D\to\overline D'$
为 $A_{K,L}$ 双模同态，
则存在次数为 $1$ 的
\[
o_{K+1,L}(\overline\varphi)
\in\operatorname{Hom}\bigl(
B_{K+1}\overline C_1[1]
\otimes\overline D[1]
\otimes B_L\overline C_0[1],
\overline D'[1]\bigr),
\]
满足下列性质：
\begin{enumerate}[label=\textup{(5.2.47.\arabic*)},leftmargin=5.8em]
\item\label{cond:5-2-47-1}
$\delta_1(o_{K+1,L}(\overline\varphi))=0$。
\item\label{cond:5-2-47-2}
作为 $\delta_1$-上同调类，
$[o_{K+1,L}(\overline\varphi)]=0$
当且仅当存在延拓
$\overline\varphi$ 的
$A_{K+1,L}$ 双模同态。
\item\label{cond:5-2-47-3}
若 $\overline\varphi'$ 与
$\overline\varphi$
是 $A_{K,L}$ 同伦的，则
\[
[o_{K+1,L}(\overline\varphi)]
=[o_{K+1,L}(\overline\varphi')]
\]
作为 $\delta_1$-上同调类。
这里 $A_{K,L}$ 同伦的概念可按显然方式定义。
\item\label{cond:5-2-47-4}
设
\[
\begin{aligned}
\overline{\mathfrak g}_{(i)}:
\overline C'_i&\longrightarrow\overline C''_i,
&
\overline{\mathfrak g}'_{(i)}:
\overline C'''_i&\longrightarrow\overline C_i
\end{aligned}
\]
为 $A_\infty$ 同态
\translatornote{原书实体页 301
把 $\overline{\mathfrak g}'_{(i)}$
的靶写作
$\overline C_{i+1}$；
但 $i=1$ 时 $C_2$ 未定义，
且下述 $\overline\varphi'':
\overline D'''\to\overline D$
位于 $\overline{\mathfrak g}'_{(i)}$
之上，右端只能对应
$\overline D$ 的基代数
$\overline C_i$。见疑误 A-0284。}。
设
$\overline\varphi':
\overline D'\to\overline D''$
和
$\overline\varphi'':
\overline D'''\to\overline D$
分别是位于
$(\overline{\mathfrak g}_{(1)},
  \overline{\mathfrak g}_{(0)})$
和
$(\overline{\mathfrak g}'_{(1)},
  \overline{\mathfrak g}'_{(0)})$
之上的 $A_{K+1,L}$ 双模同态。
则
\[
\begin{aligned}
&[o_{K+1,L}
(\overline\varphi'\circ
 \overline\varphi\circ
 \overline\varphi'')]\\
&\qquad=
(\overline\varphi'_1)_*
\circ[o_{K+1,L}(\overline\varphi)]
\circ
(\overline\varphi''_1\otimes\cdots
\otimes\overline\varphi''_1)_*.
\end{aligned}
\]
此处先保留原书的
$\overline\varphi'_1$、
$\overline\varphi''_1$ 记号
\translatornote{原书实体页 301 在
(5.2.47.4) 右端使用
$\overline\varphi'_1$ 和
$\overline\varphi''_1$，
但定义 5.2.44 对双模同态只定义
$(k_1,k_0)$ 双下标分量。
本式的线性层还涉及左右代数同态的
一次分量；原书并未说明这些单下标记号
如何同时表示两侧输入和中间双模输入。
译文保留原式、记录型别疑点，
不擅自给出未经证明的展开；见疑误 A-0285。}。
\item\label{cond:5-2-47-5}
类似地可定义
$o_{K,L+1}(\overline\varphi)$，
上列相应断言对它也成立。
\end{enumerate}
\end{lemma}

\begin{proof}
证明与定理 4.5.1 的证明相同：
置
\[
\begin{aligned}
&o_{K+1,L}(\overline\varphi)
(\mathbf x\otimes y\otimes\mathbf z)\\
&\quad=
\bigl(
\widehat{\overline d}'\circ
\widehat{\overline\varphi}_{K+1,L}
-\widehat{\overline\varphi}_{K+1,L}
\circ\widehat{\overline d}
\bigr)
(\mathbf x\otimes y\otimes\mathbf z).
\end{aligned}
\]
细节从略。
\end{proof}
